\documentclass[10pt,a4paper]{amsart}
\usepackage[margin=19mm]{geometry}
\usepackage{amsmath,amssymb,mathtools}
\usepackage[scr=rsfs,cal=euler]{mathalfa}
\usepackage{xfrac}
\usepackage[unicode,hidelinks,bookmarks=false]{hyperref}
\newcommand{\ie}{i.e.\ }
\newcommand{\cf}{cf.\ }
\newcommand{\resp}{respectively\ }
\newcommand{\Subsection}{\subsection}
\providecommand{\nobreakdash}{\nobreak\hbox{-}}
\setlength{\emergencystretch}{2em}
\multlinegap0pt
\usepackage{bm}
\usepackage{bbm}
\usepackage{dsfont}
\usepackage{xspace}
\usepackage{cancel}
\usepackage{mathtools}
\usepackage{amsthm}
\makeatletter
\@ifundefined{theorem}{\newtheorem{theorem}{Theorem}}{}
\@ifundefined{lemma}{\newtheorem{lemma}[theorem]{Lemma}}{}
\makeatother
\title[Optimal Survey Design for Private Mean Estimation]{Optimal Survey Design for Private Mean Estimation}
\author{Yu-Wei Chen, Raghu Pasupathy, Jordan A. Awan}
\date{Source: arXiv:2501.18121v1 (CC BY 4.0)}
\begin{document}
\maketitle

\noindent\emph{Source and licence.} Optimal Survey Design for Private Mean Estimation. Version arXiv:2501.18121v1 (CC BY 4.0), distributed under the Creative Commons Attribution 4.0 International licence, as recorded in the arXiv licence metadata for this exact version and shown on the abstract page https://arxiv.org/abs/2501.18121v1. Full rights notice: licenses/arxiv-2501.18121-CC-BY-4.0.txt.\par\smallskip

\noindent\textit{Editorial scope.} Counted unit: the Lemma whose source label is \texttt{lemma: change of variable y > 0} in the article (source0.tex lines 637--642, source label \texttt{lemma: change of variable y > 0}), delivered together with the article's own complete original proof of it (source0.tex lines 644--663, one \texttt{proof} environment of 20 lines). No mathematics is written, repaired or re-derived: statement and proof are the article's own text line for line. Required context retained uncounted: none. Every definition, display and label of those lines is retained unchanged. Omitted from this package and not redistributed: 1--636, 643--643, 664--785. Nothing inside a retained excerpt is omitted or altered: every retained body is delivered exactly as the article's lines are written, so no omission receipt of any kind accompanies this package. Citations: the retained excerpts carry no citation command at all, so no bibliography entry of the article is transcribed and no reference is redistributed. Reference adaptation: none was needed and none was made; every reference inside the delivered unit resolves inside it, so no unresolved reference or citation is produced. Printed slips kept literal: no printed word, symbol, hypothesis or number of the counted statement or of its proof was altered. Layout and class: the article's own class is \texttt{article} with packages including arxiv, inputenc, fontenc, hyperref, url, booktabs, siunitx, amsfonts, amsthm, amsmath, mathtools, nicefrac, microtype, graphicx; the wrapper is the lane's pinned \texttt{amsart} editorial wrapper at 10pt on A4, which restates the theorem-like environments the retained text needs and adds the article's own macro definitions that the retained text needs (none), with extra wrapper packages (bm, bbm, dsfont, xspace, cancel, mathtools). The delivered numbering is the unit's own. Assets: no retained span contains a figure environment, a graphics-inclusion command, a PDF-page inclusion, a table or a TikZ picture; no image file is copied into this package, the package declares no asset dependency and no image file is redistributed with it. Licence: version arXiv:2501.18121v1 of this article is distributed under the Creative Commons Attribution 4.0 International licence, as recorded in the arXiv licence metadata for this exact version and shown on the abstract page https://arxiv.org/abs/2501.18121v1. A full rights notice is delivered at licenses/arxiv-2501.18121-CC-BY-4.0.txt. Characters: every retained excerpt is pure ASCII, so the delivered engine, XeLaTeX with its default Latin Modern text font, reports no missing character.
\begin{lemma} \label{lemma: change of variable y > 0}
    For all $y>0$,
    \begin{align} \label{eq: change of variable y}
    \left(\left(1+\frac{1}{y}\right) \log(1+y) - 2\right)^2 - 1 + \log(1+y) > 0.
    \end{align}
\end{lemma}

\begin{proof}
    %We prove in two ways that \eqref{eq: change of variable y} $>0$ for all $y>0$. One way is to consider derivatives. 
    First, since we focus on the positive real line ($y > 0$), we can do change of variable by letting $y = e^x - 1$ with domain $x > 0$, then \eqref{eq: change of variable y} becomes
    \begin{align*}
        \left( \frac{x e^x}{e^x - 1} - 2 \right)^2 - 1 + x
        .
    \end{align*}
    We can multiply through by $(e^x -1)^2$ to give us the equivalent of problem of establishing that $h(x)>0$, where
    \begin{align*}
        h(x) \coloneqq (x e^x - 2 e^x + 2)^2 + (x-1) (e^x - 1)^2.
    \end{align*}
    To analyze $h(x)$, we compute its derivatives up to the third order:
    \begin{align*}
    \frac{d h(x)}{d x} =& e^{2x}\left(2x^{2}-4x+3\right)+2e^{x}\left(x-2\right)+1 , \\
    \frac{d^2 h(x)}{d x^2} =& 2e^x\left(e^x\left(2x^2-2x+1\right)+x-1\right) , \\
    \frac{d^3 h(x)}{d x^3} =& 2 x e^x (4 x e^x + 1).
    \end{align*}

    We begin by noting that \( \frac{d^3 h(x)}{d x^3} > 0 \) for all \( x > 0 \), which implies that \( \frac{d^2 h(x)}{d x^2} \) is increasing. Since \( \frac{d^2 h(x)}{d x^2} \big|_{x=0} = 0 \), it follows that \( \frac{d^2 h(x)}{d x^2} > 0 \) for all \( x > 0 \). By the same argument, knowing that \( \frac{d h(x)}{d x} \big|_{x=0} = 0 \), we conclude that \( \frac{d h}{d x} > 0 \) for all \( x > 0 \). Finally, since \( h(0) = 0 \), we establish that \( h(x) > 0 \) for all \( x > 0 \).
\end{proof}

\end{document}
